\chapter{数据准备：从原始记录到可用特征}
\label{ch:data}

模型接收的是经过整理的输入。列代表什么、缺失怎样表达、尺度如何确定，
都会改变模型能学到的关系。本章用小表和可复算的数字解释这些步骤。

\begin{keyidea}[先划分，再学习处理规则]
凡是需要从数据中估计的处理规则，都属于拟合流程的一部分。
填补值、缩放参数、类别字典和特征选择规则只在训练数据上学习；
验证与测试沿用这些规则。交叉验证时，这个边界要在每一折重新建立。
\end{keyidea}

\section{样本矩阵：行、列与目标}
\label{sec:data-matrix}
样本矩阵（Design Matrix）$X\in\R^{n\times d}$ 的每行是一条样本，
每列是一个输入特征。$x_{ij}$ 表示第 $i$ 个样本的第 $j$ 个特征。
单目标监督学习有长度为 $n$ 的目标向量 $\vct y$：
数值回归中常写作 $\vct y\in\R^n$，分类时每项则是一个类别标签。

\begin{figure}[htbp]
\centering\begin{tikzpicture}[x=18mm,y=9mm]
\foreach \r/\a/\b/\c in {1/10/2/0,2/20/1/1,3/30/4/0}{
 \foreach \j/\v in {1/\a,2/\b,3/\c}{
  \draw[draw=Teal!45,fill=Pale] (\j,-\r) rectangle ++(1,-1);
  \node at (\j+.5,-\r-.5) {$\v$};}
}
\node[lab] at (2.5,-.4) {$X\in\R^{3\times3}$：三列输入特征};
\node[lab,align=right] at (-.1,-2.5) {每行\\一个样本};
\draw[flow] (.65,-2.5)--(1,-2.5);
\foreach \r/\v in {1/0,2/1,3/1}{
 \draw[draw=Amber,fill=Amber!5] (5,-\r) rectangle ++(1,-1);
 \node at (5.5,-\r-.5) {$\v$};}
\node[lab] at (5.5,-.4) {$y$：目标};
\node[lab] at (3.5,-4.7) {输入与目标按行对应，目标不作为输入列};
\end{tikzpicture}

\caption{三个样本、三个数值特征与一个目标的结构示例。数字仅用于说明形状。}
\par\smallskip
\readfigure{第三列只有0和1，仍可能是输入特征。是不是目标取决于任务定义，
不能仅凭数据类型或列的位置判断。样本顺序改变时，$X$与$y$必须同步改变。}
\end{figure}

\subsection{先问预测时能否获得这个特征}
预测违约时，申请时的收入可能可用，事后产生的催收记录不能提前作为输入。
姓名、设备编号、时间戳不应机械地当作普通连续数值：它们可能用于分组划分、
构造合理特征，也可能让模型记住身份或未来信息。

\quick{特征不是“表中除标签之外的所有列”；每一列都需要含义、单位和可用时间。}

\section{无序类别与独热编码}
\label{sec:data-onehot}
无序类别（Nominal Category）只有类别差异，没有自然大小顺序。
把红、蓝、绿编码为0、1、2，并不能说明绿比蓝“大”，也不能说明两者距离为1。
独热编码（One-Hot Encoding）为每个类别建立一个指示列。

\begin{figure}[htbp]
\centering\begin{tikzpicture}[x=22mm,y=9mm]
\node[lab] at (0,0) {原类别};
\foreach \r/\c in {1/红,2/蓝,3/绿}{\node at (0,-\r) {\c};}
\draw[flow] (.5,-2)--(1.2,-2);
\foreach \j/\c in {2/红,3/蓝,4/绿}{\node[lab] at (\j,0) {\c};}
\foreach \r in {1,2,3}{\foreach \j in {2,3,4}{
 \pgfmathtruncatemacro{\v}{\j==\r+1?1:0}
 \node[draw=Teal!30,fill=Pale,minimum width=20mm,minimum height=8mm] at (\j,-\r) {$\v$};}}
\node[lab] at (2,-4) {三列分别回答“是不是该类别”};
\end{tikzpicture}

\caption{三种颜色映射到三个指示列；这里没有定义颜色间的顺序。}
\par\smallskip
\readfigure{一条非缺失、属于已知类别的记录，在完整独热编码中恰有一列为1。
拟合时记录类别字典和列顺序；预测时必须保持同一顺序。}
\end{figure}

\subsection{三个常见边界}
训练中未出现的新类别需要预先定义处理方法，例如报错、归入“其他”，或按编码器约定输出全零。
若全零同时用作基准类别，应理解其语义可能重合。
类别很多时，独热矩阵会很宽；可考虑稀疏存储、合理合并或专门的类别模型。
有些树模型原生支持类别特征，应遵循实现接口。

\begin{pitfall}[独热列与截距可能线性相关]
完整的$K$个独热列逐行之和为1，和截距列相关。
无约束线性模型常删除一个基准类别列，或采用相应约束；
是否删除要结合模型与正则化，不能把“所有模型必须删一列”当成规则。
\end{pitfall}

\section{有序编码保留顺序，不保证间距}
\label{sec:data-ordinal}
有序类别（Ordinal Category）具有自然顺序，例如低、中、高。
编码为0、1、2可以表达这种顺序，但不保证“中到高”与“低到中”的实际差异相等。

\begin{center}
\begin{tabularx}{\textwidth}{P{27mm}P{30mm}Y}
\toprule
原特征 & 可用表示 & 引入的含义\\
\midrule
低、中、高 & 有序编码0、1、2 & 保留顺序；在线性模型中还引入等间距效应。\\
红、蓝、绿 & 三个独热列 & 各类别分别具有指示特征，不强加大小关系。\\
真实温度 & 有单位的数值 & 数值差异有物理意义，需保留测量单位。\\
\bottomrule
\end{tabularx}
\end{center}

例如线性预测 $\hat y=b+wx$，低、中、高编码0、1、2后，相邻等级的预测差都为$w$。
这是一项模型假设；若不合理，可以用独热编码分别估计各等级效果。
树对有序整数作阈值切分可以保留等级顺序，但任意给无序类别编号会限制它可直接表达的分组。

\begin{pitfall}[编码值不是比例尺度]
低=1、中=2、高=3，只说明人为选定了三个数字。
不能据此说“高的实际程度是低的三倍”。顺序、差值与倍数是不同层次的信息。
\end{pitfall}
\quick{编码方式要与变量含义及模型的使用方式一起选择。}

\section{缺失值：不知道，不等于零}
\label{sec:data-missing}
缺失值（Missing Value）表示未获得观测。零可能是有效数值，不能默认用零代替所有缺失。
先区分未测量、不适用、采集失败等原因，再选择处理策略。

\begin{center}
\begin{tabularx}{\textwidth}{P{32mm}Y}
\toprule
处理方式 & 适用考虑与代价\\
\midrule
训练均值或中位数 & 数值特征的简单基线；压缩变异，不能恢复真实缺失值。\\
训练众数或专门类别 & 用于类别列；专门类别可以显式保留“未知”的语义。\\
增加缺失指示列 & 让模型知道该位置原本缺失；需防止采集流程变化破坏关系。\\
删除行或列 & 损失信息，并可能改变样本分布；不能只凭缺失比例小就断言无偏。\\
利用其他特征预测 & 如近邻填补；依赖相关结构，仍需验证效果。\\
\bottomrule
\end{tabularx}
\end{center}

训练列为$[10,\text{缺失},20,30]$时，中位数为20。
填补后为$[10,20,20,30]$；缺失指示列为$[0,1,0,0]$。
新样本的缺失仍用训练得到的20填补，不在测试列上重新计算中位数。

\begin{pitfall}[时间顺序也约束填补]
预测某时刻的结果时，不能利用当时尚未产生的未来记录反向填补。
有些模型可以直接处理缺失值，也需要确认实现的支持范围和训练、预测时的语义一致。
\end{pitfall}
\quick{填补后的值是建模用的估计，不是已经找回的真实观测。}

\section{近邻填补：用观测部分寻找相似记录}
\label{sec:data-knn}
$k$近邻填补（$k$-NN Imputation）根据已观测特征寻找相似的训练记录，
再汇总这些邻居在缺失列上的值。下面只用一个已观测特征$A$计算距离。

\begin{figure}[htbp]
\centering\begin{tikzpicture}
\begin{axis}[width=133mm,height=43mm,xmin=0,xmax=10,ymin=-.3,ymax=1.6,
 xlabel={观测到的特征 $A$},ytick=\empty,axis y line=none,axis x line=bottom,
 xtick={1,2,2.1,8,9},xticklabels={1,2,,8,9}]
\addplot[only marks,mark=*,mark size=3pt,Teal] coordinates {(1,0)(2,0)(8,0)(9,0)};
\addplot[only marks,mark=diamond*,mark size=4pt,Amber] coordinates {(2.1,.6)};
\draw[flow,dashed] (axis cs:2.1,.6)--(axis cs:1,0);
\draw[flow,dashed] (axis cs:2.1,.6)--(axis cs:2,0);
\node[lab,above] at (axis cs:2.1,.8) {$A=2.1,\ B=?$};
\node[lab,above left] at (axis cs:1,.05) {$B=10$};
\node[lab,below right] at (axis cs:2,-.02) {$B=12$};
\node[lab,above] at (axis cs:8,.05) {$B=80$};
\node[lab,above] at (axis cs:9,.05) {$B=90$};
\end{axis}
\end{tikzpicture}

\caption{训练记录为$(1,10),(2,12),(8,80),(9,90)$；橙色新记录的$B$缺失。}
\par\smallskip
\readfigure{到$A=2.1$的距离分别为1.1、0.1、5.9、6.9。
当$k=2$且等权时，选择前两条记录，填补$B=(10+12)/2=11$。
全部训练$B$的中位数则为$(12+80)/2=46$；两种规则表达不同假设。}
\end{figure}

\subsection{距离怎样定义，决定谁是邻居}
缺失的$B$本身不能用于寻找邻居。多维特征可依据共同观测的部分定义距离，
但共同观测很少时相似度可能不可靠；具体距离和无可用邻居时的回退规则需明确。
量纲也会影响距离，例如以元计价的收入可能淹没以年计的年龄。

\begin{pitfall}[没有固定最优的邻居数]
$k$、权重、尺度与距离都应通过训练内验证选择。
仅从训练数据保存供填补使用的邻居，不能把验证或测试记录加入邻居库。
近邻填补可能保留局部结构，也可能传播异常值，不能预先断言优于简单填补。
\end{pitfall}

\section{缩放：参数来自训练集}
\label{sec:data-scaling}
Min--Max缩放使用训练最小值$a$和最大值$b$：
\[
 x'=\frac{x-a}{b-a}\quad (b>a).
\]
标准化（Standardization）使用训练均值$\mu$与标准差$s$：
\[
 x'=\frac{x-\mu}{s},\qquad
 \mu=\frac1n\sum_i x_i,\quad s=\sqrt{\frac1n\sum_i(x_i-\mu)^2}.
\]
这里明确使用分母$n$，与一些统计教材的样本标准差约定不同。

\begin{figure}[htbp]
\centering\begin{tikzpicture}[x=26mm,y=12mm]
\foreach \r/\name in {0/原值,1/Min--Max,2/标准化}{
 \node[lab,anchor=east] at (-.35,-\r) {\name};
 \draw[->,draw=Quiet] (-.1,-\r)--(3.45,-\r);}
\foreach \j/\raw/\mm/\std in {0/10/0/{-1.225},1/20/{0.5}/0,2/30/1/{1.225},3/40/{1.5}/{2.449}}{
 \foreach \r/\v in {0/\raw,1/\mm,2/\std}{
  \ifnum\j=3\relax\def\cc{Amber}\else\def\cc{Teal}\fi
  \fill[\cc] (\j,-\r) circle (2pt);
  \node[font=\small,above,text=\cc] at (\j,-\r) {$\v$};}}
\draw[decorate,decoration={brace,mirror,amplitude=4pt},Teal] (-.15,-2.35)--(2.15,-2.35);
\node[lab] at (1,-2.85) {训练数据：拟合缩放参数};
\node[lab,align=center,text=Amber] at (3,-2.8) {新样本\\沿用参数};
\end{tikzpicture}

\caption{训练值10、20、30；新值40沿用同一套参数。标准化数字保留三位小数。}
\par\smallskip
\readfigure{训练均值为20，标准差为$\sqrt{200/3}\approx8.165$。
新值40经Min--Max变为1.5，经标准化变为约2.449。
训练范围映射到$[0,1]$，不意味着所有未来值都在$[0,1]$内。}
\end{figure}

\begin{pitfall}[标准化不会自动变成正态分布]
减均值、除标准差只改变位置与尺度。它不消除偏态，也不消除异常值。
常数列分母为零，需要明确保留、删除或特殊处理规则，不能直接相除。
\end{pitfall}
距离、正则化和梯度更新常受尺度影响；常规按单列阈值切分的树通常不依赖这种缩放。
是否缩放仍由模型与特征含义决定。

\section{特征选择：筛选规则也是模型的一部分}
\label{sec:data-selection}
特征选择（Feature Selection）保留原特征的一个子集。
它不同于把原特征重新组合成新坐标的特征提取，例如主成分分析。

\begin{center}
\begin{tabularx}{\textwidth}{P{26mm}P{49mm}Y}
\toprule
方法 & 怎样选择 & 主要边界\\
\midrule
过滤法 Filter & 用统计量或单变量分数筛选。 & 快；单变量分数可能漏掉交互作用。\\
包装法 Wrapper & 对候选子集拟合并验证模型，如前向选择、递归消除。 & 开销大；反复比较会过拟合验证结果。\\
嵌入法 Embedded & 在训练中形成选择，如L1惩罚产生零系数。 & 依赖模型、尺度及超参数。\\
\bottomrule
\end{tabularx}
\end{center}

先在全部数据上用标签选出“相关性最高”的列，再进行交叉验证，已经泄漏验证标签。
每折需要只用子训练集拟合筛选器，再把选定列应用到该折验证集。

\subsection{不要把经验规则写成定理}
高度相关的列可能冗余，也可能各自包含有效信息；低方差的指示特征仍可能识别重要的罕见类别。
树的重要性、线性系数大小与预测贡献都受模型和数据影响，不能直接解释为因果作用。

\quick{先确定选择流程，再评估整个流程。删列后的泛化表现，比“看起来显著”更接近实际目标。}

\section{类别不平衡：调整训练，不改写评估对象}
\label{sec:data-sampling}
类别不平衡（Class Imbalance）指类别频率差异较大。
多数类、少数类由频率决定，正类、负类由任务定义决定，两组概念不能混用。

\begin{figure}[htbp]
\centering\begin{tikzpicture}[x=1mm,y=1mm]
\foreach \y/\major/\minor/\name in {0/90/10/原训练集,-18/20/10/欠采样示例,-36/90/45/过采样示例}{
 \node[lab,anchor=east] at (0,\y+3) {\name};
 \fill[Teal!70] (3,\y) rectangle ++(\major*.65,6);
 \fill[Amber!75] (3+\major*.65,\y) rectangle ++(\minor*.65,6);
 \node[lab,anchor=west] at ({7+(\major+\minor)*.65},\y+3) {\major\,负 / \minor\,正};}
\node[lab,anchor=west] at (3,-48) {青色：多数类\qquad 橙色：少数类};
\end{tikzpicture}

\caption{90个负类、10个正类的两个独立采样示例；不要求一定变成1:1。}
\par\smallskip
\readfigure{欠采样（Undersampling）从多数类中保留较小子集，会损失部分记录。
随机过采样（Random Oversampling）重复抽取少数类，使其在训练中出现更多次；
45条少数类训练记录不代表获得45个独立的新对象。}
\end{figure}

采样只在训练数据内进行；验证与测试保留目标使用场景的分布。
交叉验证时也只处理每折子训练集。先重复样本再划分，会使相同对象跨训练与验证。

\begin{pitfall}[比例变化会影响概率解释]
采样改变了模型接触的类别先验。输出的概率未必仍是原始人群中的风险概率。
阈值选择与概率校准需要在具有目标分布、用途独立的数据上处理。
也可以比较类别权重等方案，不必先认定采样一定有益。
\end{pitfall}

\section{SMOTE：在少数类邻居之间插值}
\label{sec:data-smote}
SMOTE是一种合成过采样方法。基本形式在少数类样本$x_i$与其少数类近邻$x_j$之间生成
\[
 x_{\mathrm{new}}=x_i+\lambda(x_j-x_i),\qquad0\leq\lambda\leq1.
\]
若$x_i=(1,2)$、$x_j=(3,4)$且$\lambda=0.25$，则$x_{\mathrm{new}}=(1.5,2.5)$。
这条新记录位于两点之间，并不证明它在实际世界中一定合理。

\begin{center}
\begin{tabularx}{\textwidth}{P{37mm}Y}
\toprule
需要检查 & 原因\\
\midrule
特征尺度与距离 & 它们决定选到哪些近邻。\\
少数类边界与异常值 & 插值可能穿过多数类区域，或放大噪声。\\
类别型与离散约束 & 直接插值可能产生无意义的“半个类别”；需相应方法或表示。\\
少数类数量 & 太少时近邻关系不稳，且可能无法满足算法的邻居数要求。\\
独立验证表现 & 合成点不是独立采集的额外证据，仍需在目标分布上评估。\\
\bottomrule
\end{tabularx}
\end{center}

\quick{随机复制改变记录出现频率；SMOTE改变训练点的位置。两者都不能替代真实的新数据。}

\section{把处理步骤放进同一条流水线}
\label{sec:data-pipeline}
流水线（Pipeline）把处理步骤按顺序连接，让训练和预测遵循同一套规则。
数值列与类别列可以分别处理，再拼接为模型输入。

\begin{figure}[htbp]
\centering\begin{tikzpicture}[node distance=7mm and 8mm]
\node[box,text width=20mm] (num) {数值列};
\node[box,text width=23mm,right=of num] (imp) {填补缺失};
\node[box,text width=23mm,right=of imp] (scale) {缩放};
\node[box,text width=20mm,below=12mm of num] (cat) {类别列};
\node[box,text width=23mm,right=of cat] (fill) {缺失类别};
\node[box,text width=23mm,right=of fill] (hot) {独热编码};
\node[box,text width=20mm,at={($(scale)!0.5!(hot)+(3.25,0)$)}] (join) {拼接特征};
\node[box,text width=20mm,right=of join] (model) {模型};
\draw[flow] (num)--(imp);\draw[flow] (imp)--(scale);
\draw[flow] (cat)--(fill);\draw[flow] (fill)--(hot);
\draw[flow] (scale.east)-|(join.north);\draw[flow] (hot.east)-|(join.south);
\draw[flow] (join)--(model);
\end{tikzpicture}

\caption{一个基础表格数据流程；箭头表示数据经过的步骤。可在训练支路添加筛选和采样。}
\par\smallskip
\readfigure{拟合时，在训练集学习各步所需参数并拟合模型。
预测时只调用已学规则。数值填补放在缩放之前，是此例的选择；
近邻填补等流程还需单独设计距离尺度，不能机械套用。}
\end{figure}

\begin{center}
\begin{tabularx}{\textwidth}{P{32mm}YY}
\toprule
步骤 & 训练时学习 & 验证、测试或上线时使用\\
\midrule
填补 & 中位数、众数或邻居库 & 固定填补规则\\
缩放 & 均值、标准差或极值 & 固定参数变换\\
编码 & 类别字典与列顺序 & 同一字典及未知类策略\\
选择 & 训练内确定的特征子集 & 同一组列\\
采样 & 训练记录的选择、复制或合成 & 不对待评估记录采样\\
\bottomrule
\end{tabularx}
\end{center}

\subsection{自测与回看}
测试值超出训练最大值，经Min--Max得到1.2，是否一定是程序错误？\textbf{不是。}
交叉验证前用全部标签筛选列，是否合适？\textbf{不合适，筛选应在每折训练内。}
复制10条少数类到50条，是否等于新增40个独立对象？\textbf{不等于。}

\smallskip
\begingroup\small
\noindent\textbf{核查依据：}scikit-learn的Preprocessing、Imputation、Feature selection与Common pitfalls文档，
以及imbalanced-learn的Common pitfalls文档（访问于2026-10-28）。
填补、缩放与采样的数值例子可用配套源码中的\texttt{verify\_data.py}复核。
\par\endgroup
